%================================
% Vocabulary & Language Focus — XB1 (A–C)
% 编译方式: xelatex vocab-handout.tex (需运行两遍以生成目录)
%================================
\documentclass[../main.tex]{subfiles}

\begin{document}

%================================
\section{A}
%================================

\vw{abandon}{/əˈbændən/}{vt.}

\sense{vt.}{中止，放弃，不再有}
\col{abandon sth.}
\ex{abandon a plan / an idea / an attempt}
\ex{They had to abandon the match because of the rain.}
\ex{The original design had to be abandoned when the budget was cut.}
\col{abandon hope (of doing sth.)}
\ex{Rescuers have not abandoned hope of finding more survivors.}

\sense{vt.}{抛弃（人或动物）；舍弃（地方）}
\col{abandon sb.}
\ex{He felt he had been abandoned by his family.}
\col{abandon oneself to sth.（沉溺于，书面语）}
\ex{They abandoned themselves to grief when they heard the news.}

\sense{adj.}{abandoned 废弃的}
\ex{an abandoned car / building}
\ex{The village was abandoned after the flood of 1954.}

\begin{origin}
源自古法语短语 \textit{(mettre) a bandon}“（把某物）交由某人处置”，其中 \textit{bandon} 意为“权力、管辖”，来自日耳曼语词根 *\textit{ban}“宣告、命令”（同源词 \textbf{ban} 禁令）。字面义即“任凭他人意志摆布”，引申为“完全撒手不管”$\to$“放弃”。这是典型的“短语凝固成单词”的构词方式，类似的还有 \textit{alone}（all + one）。
\end{origin}

\vw{accompany}{/əˈkʌmpəni/}{vt.}

\sense{vt.}{陪同，陪伴}
\col{accompany sb. to spl.（注意：介词用 to，不加 with）}
\ex{Children under 12 must be accompanied by an adult.}
\ex{The President was accompanied by his wife on the visit.}

\sense{vt.}{伴随，与……同时发生}
\col{be accompanied by/with sth.}
\ex{Strong winds were accompanied by heavy rain throughout the night.}

\sense{vt.}{为……伴奏}
\ex{She sang the folk song, accompanied by her brother on the guitar.}

\begin{origin}
经由法语 \textit{accompagnier} 进入英语，由 à- + \textit{compaignon}（同伴）构成；而 \textit{compaignon} 追溯至晚期拉丁语 \textit{companionem}，字面义为“\textbf{分面包的人}”：\textit{com-}（共同）+ \textit{panis}（面包，源自原始印欧语 *\textit{pa-}“喂养”）。“同伴”的本义竟是“与你掰同一块面包的人”——正如汉语“伙计”的“伙”本指共火而食。
\end{origin}

\begin{gram}[accompany 不需要 with]
汉语说“陪同某人”，英语直接说 \textit{accompany sb..}。\textcolor{WildStrawberry!80!black}{\ding{55}} \textit{accompany with sb..} 是典型中式错误。表“陪同去某地”用 \textit{accompany sb.. \textbf{to}} a place。名词形式：\textbf{company}（陪伴；for company 作伴）、\textbf{companion}（同伴）。
\end{gram}

\vw{adaptation}{/ˌædæpˈteɪʃn/}{n.}

\sense{n.}{改编本，改写本}
\col{an adaptation of sth.}
\ex{The film is an adaptation of a best-selling novel.}
\ex{the stage adaptation of \textit{Hamlet}}

\sense{n.}{适应}
\col{adaptation to sth.}
\ex{The camel's hump is a remarkable adaptation to desert life.}

\begin{lex}{adapt / adopt / adjust}
\textbf{adapt} $v.$（使适应；改编）：\textit{adapt (oneself) to the new environment}，名词即 adaptation；\textbf{adopt} $v.$ 收养；采纳（拼写仅一字母之差，常考）；\textbf{adjust} $v.$ 调节（多指机械性微调）：\textit{adjust the seat height}。
\end{lex}

\begin{origin}
拉丁语 \textit{adaptare}“使贴合”：\textit{ad-}（朝向）+ \textit{aptare}“使适合”，来自 \textit{aptus}“贴合的、恰当的”（即 apt 一词：\textit{an apt description} 恰当的描述）。“改编”与“适应”共享同一意象——把某物\textbf{贴合}到新环境或新载体上。
\end{origin}

\vw{adorable}{/əˈdɔːrəbl/}{adj.}

\sense{adj.}{可爱的，讨人喜爱的}
\ex{an adorable baby / puppy / smile}
\ex{The kitten looked adorable chasing its own tail.}

\begin{origin}
拉丁语 \textit{adorari}“向……祈祷、恳求”：\textit{ad-} + \textit{orare}“说话、祈祷”。同源词 \textbf{oracle}（神谕，预言者之口）、\textbf{orator}（演说家）、\textbf{oral}（口头的）。宗教语境中 adore 本指“向神祷告、礼赞”，14 世纪后才软化出“万分喜爱”之义——“可爱得让人想拜一拜”。
\end{origin}

\vw{advocate}{/ˈædvəkeɪt/ vt.\ \ /ˈædvəkət/ n.}{vt. \& n.}

\sense{vt. /ˈædvəkeɪt/}{拥护，支持，提倡}
\col{advocate (doing) sth.（注意：接动名词，不接不定式）}
\ex{Many experts advocate reducing the amount of sugar in children's diets.}
\ex{The group advocates tighter gun control.}

\sense{n. /ˈædvəkət/}{拥护者，支持者；辩护律师}
\col{an advocate of/for sth.}
\ex{She is a passionate advocate of animal rights.}
\ex{The court appointed an advocate for the defendant.}

\begin{gram}[名动异读：重音转移构词法]
英语中一批双音节词，\textbf{动词重音在后、名词重音在前}，且尾元音弱化：\textit{adv\textbf{o}cate}（名，/-kət/）vs.\ \textit{advoc\textbf{a}te}（动，/-keɪt/）。同类词对：
\begin{center}
\begin{tabular}{llll}
\toprule
动词 & 名词 & 动词 & 名词 \\
\midrule
re\textbf{cord} 记录 & \textbf{re}cord 唱片；记录 & per\textbf{mit} 允许 & \textbf{per}mit 许可证 \\
im\textbf{port} 进口 & \textbf{im}port 进口货 & pro\textbf{duce} 生产 & \textbf{pro}duce 农产品 \\
con\textbf{trast} 对比 & \textbf{con}trast 对比 & in\textbf{crease} 增加 & \textbf{in}crease 增加 \\
\bottomrule
\end{tabular}
\end{center}
判断词性的快捷方法：重音在第二音节 $\Rightarrow$ 多为动词。
\end{gram}

\begin{origin}
拉丁语 \textit{advocatus}“\textbf{被传召来帮忙的人}”：\textit{ad-}（向）+ \textit{vocatus}（\textit{vocare}“呼唤”的过去分词）。词根 \textit{voc-/vox}“呼唤、声音”来自原始印欧语 *\textit{wekw-}“说话”，同族词有 \textbf{voice}、\textbf{vowel}（元音）、\textbf{voca\-tive}（呼格）、\textbf{convoke}（召集）、\textbf{vocation}（天职、神召）。“辩护律师”正是“被喊来替你说话的人”。
\end{origin}

\vp{all of a sudden}

\sense{adv. 短语}{突然，猛地}
\ex{All of a sudden, the lights went out and the room was in total darkness.}

\begin{gram}[名词短语作状语]
“all of a sudden” 中名词 sudden 前用了不定冠词 a，这是“名词短语作状语”的用法，类似的还有 \textit{in a hurry}（匆忙地）、\textit{with a smile}（微笑着）、\textit{all at once}（一下子）。写作中用 all of a sudden 开头能给叙述制造突转效果。
\end{gram}

\vp{all walks of life}

\sense{n. 短语}{各行各业，社会各界}
\col{people from all walks of life（固定介词 from）}
\ex{The charity attracts volunteers from all walks of life --- doctors, teachers, truck drivers and students.}
\ex{The festival is designed to appeal to people of all walks of life.}

\begin{origin}
“walk of life” 中的 walk 保留其古义“\textbf{行业、营生}”（与德语 \textit{Wandel}“生涯”同源）。这个短语由“神职人员、贵族、平民”等人生“路径”的隐喻而来，18 世纪定型。注意 walk 在此处不可按“行走”理解——这是“历史语义残留”的典型例子。
\end{origin}

\vw{allowance}{/əˈlaʊəns/}{n.}

\sense{n.}{津贴，补助；零花钱}
\ex{a clothing / travel / housing allowance}
\ex{Employees receive a monthly travel allowance of 300 yuan.}
\col{give sb. an allowance（美式；英式多说 pocket money）}
\ex{My parents gave me a small weekly allowance for doing housework.}

\sense{n.}{限额}
\col{a baggage allowance of 23 kg}
\ex{Each passenger is entitled to a baggage allowance of 23 kg on this flight.}

\vp{make allowance(s) for}

\sense{v. 短语}{体谅；考虑到，估计到}
\ex{You must make allowances for him --- he is only six and it's his first performance.}
\ex{The schedule makes allowance for possible delays caused by bad weather.}

\begin{origin}
allowance 与 allow 同源，经由古法语 \textit{alouer} 追溯至拉丁语 \textit{allocare}“\textbf{拨给、分配位置}”：\textit{ad-} + \textit{locare}“放置”（\textit{locus}“地方”，同源词 \textbf{location}、\textbf{local}）。allow 的语义链是“分派 $\to$ 批准 $\to$ 允许”，而 allowance 在 15 世纪先表“承认、认可”，后经“核定拨给的份额”才有了“津贴”义—— allowances 是“被\textbf{划拨}出来的钱”，不是“被允许的钱”。
\end{origin}

\vw{ambitious}{/æmˈbɪʃəs/}{adj.}

\sense{adj.}{宏大的，艰巨的（修饰事物）}
\ex{an ambitious plan / project / target}
\ex{Building the canal in just five years was an extremely ambitious undertaking.}

\sense{adj.}{有野心的，有雄心的（修饰人）}
\col{be ambitious for sth. / to do sth.}
\ex{He is ambitious to become a professional violinist.}
\ex{Parents are often too ambitious for their children.}

\sense{n.}{ambition 野心，雄心}
\col{realize / achieve one's ambition}
\ex{He finally achieved his burning ambition of climbing Everest.}

\begin{origin}
拉丁语 \textit{ambitio}“\textbf{四处奔走（拉票）}”，来自动词 \textit{ambire}：\textit{ambi-}“环绕、两边”（同源词 \textbf{ambulance}“流动医院”，本义“边走边漂移的野战医院”、\textbf{ambiguous}“歧义的，字面‘两边都走’”）+ \textit{ire}“走”（不定式 \textit{ire}，同源词 \textbf{exit}“他出去”、\textbf{circuit}“绕行”）。古罗马候选人求官须\textbf{绕着}广场挨个握手拉票，“雄心”的本义竟是“竞选时四处兜圈”。
\end{origin}

\vw{apart}{/əˈpɑːt/}{adv.}

\sense{adv.}{成碎片；分开；相隔}
\col{take / pull / tear sth. apart}
\ex{He took the clock apart to see how it worked, but couldn't put it back together.}
\col{(live) apart / apart from each other}
\ex{The twins were brought up apart and only met again after twenty years.}
\col{set sb../sth. apart (from\dots)（使与众不同）}
\ex{Her gift for languages set her apart from her classmates.}
\col{数字 + apart（相隔多少距离、时间、分数）}
\ex{The two elections were only three months apart.}
\col{fall apart（破碎，破裂；崩溃）}
\ex{The old book literally fell apart in my hands.}
\ex{After his wife died, the family began to fall apart.}
\col{apart from（兼有 besides 与 except for 两义）}
\ex{Apart from a few s[l.elling mistakes, your essay is excellent.（= except for）}
\ex{Apart from English, she also s[l.eaks French.（= besides）}

\begin{origin}
\textit{a-}（处于，古英语 \textit{on} 的弱化形式）+ \textit{part}“部分”，字面即“成部分地（分开）”。\textit{part} 来自拉丁语 \textit{pars, partis}“一份、份额”（原始印欧语 *\textit{perə-}“授予、分配”）。同结构的副词 \textbf{aside}（a- + side）、\textbf{aboard}（a- + board 甲板）、\textbf{alive}（a- + life）中的 a- 都是同一个“在……之上/处于”的前缀。
\end{origin}

\vw{apparent}{/əˈpærənt/}{adj.}

\sense{adj.}{显而易见，明白易懂，显然}
\col{It is apparent (to sb..) that\dots}
\ex{It was apparent to everyone that the company was in serious trouble.}
\col{become apparent}
\ex{It soon became apparent why he had refused the offer.}

\begin{lex}
apparent 既表“显然的”（=\,evident），也表“\textbf{表面的、貌似的}”（此义为它独有）：\textit{His apparent calm concealed great anxiety.} 他表面的平静掩藏着极大的焦虑。阅读理解常在“表面的”这一义项上设题，切勿见 apparent 即译“明显的”。辨析：\textbf{obvious} 语气最强，指一目了然；\textbf{evident} 较书面，指有证据支撑；\textbf{apparent} 兼有上述两义，须依语境判断。
\end{lex}

\vw{appetite}{/ˈæpɪtaɪt/}{n.}

\sense{n.}{食欲，胃口}
\col{have a good / poor appetite}
\ex{The fresh air and exercise gave us a good appetite.}
\col{lose / spoil / ruin one's appetite}
\ex{Don't eat sweets before dinner --- you'll spoil your appetite.}

\sense{n.}{强烈欲望}
\col{an appetite for sth.（固定介词 for）}
\ex{The public has an insatiable appetite for celebrity gossip.}

\sense{adj.}{appetizing 开胃的，诱人的}
\ex{The smell of fresh bread is very appetizing.}

\begin{origin}
拉丁语 \textit{appetitus}“向……伸手的渴望”，动词 \textit{appetere}：\textit{ad-}（朝向）+ \textit{petere}“\textbf{追寻、请求、冲击}”（原始印欧语 *\textit{pet-}“冲、扑向”）。同根词构成一大“追求”家族：\textbf{compete}（com- 一起追求 $\to$ 竞争）、\textbf{appetite} 与 \textbf{impetus}（冲击力）、\textbf{perpetual}（per- 一直追求下去 $\to$ 永恒的）、\textbf{petition}（请愿书）、\textbf{repeat}（re- 再次求取 $\to$ 重复）。哲学史上的“ appetitive soul”（欲望灵魂，与理性灵魂相对）即用此词。
\end{origin}

\vw{appropriate}{/əˈprəʊpriət/}{adj.}

\sense{adj.}{合适的}
\col{appropriate (for sb../sth.) / appropriate to do sth.}
\ex{Jeans are not appropriate for a formal interview.}
\ex{It would not be appropriate to discuss the matter now.}
\col{it is appropriate that\dots（从句常用 should 型虚拟语气）}
\ex{It is appropriate that a special committee (should) be set up to look into the issue.}

\begin{origin}
拉丁语 \textit{appropriare}“\textbf{收归己有}”：\textit{ad-} + \textit{proprius}“自己的、特有的”。\textit{proprius} 一般析为 \textit{pro privo}“为个人而设”（\textit{privus}“单个的”，同源词 \textbf{private}、\textbf{deprive}“夺走”）。故 appropriate 作形容词即“归属性恰当——恰好属于此情此景”；作动词则仍是“挪作己用”（\textit{appropriated funds} 挪用的款项）。同根词 \textbf{property}（财产；特性“专属属性”）、\textbf{proper}（恰当的；专有的）。
\end{origin}

\vw{architecture}{/ˈɑːkɪtektʃər/}{n.}

\sense{n.}{建筑设计，建筑风格}
\ex{The city is famous for its colonial-era architecture.}

\sense{n.}{建筑学}
\ex{She studied architecture at Tsinghua University.}
\col{the architecture of\dots（引申：（系统的）架构）}
\ex{the architecture of the new operating system}
\col{the architect of sth.（建筑师；喻“缔造者”）}
\ex{the chief architect of the economic reform}

\begin{origin}
希腊语 \textit{arkhitektōn}“\textbf{工头、总建造师}”：\textit{arkhi-}“为首的”（\textit{arkhein}“统治、居首”，同源词 \textbf{archbishop} 大主教、\textbf{monarch} 君主、\textbf{anarchy} 无统治$\to$无政府状态）+ \textit{tekton}“建造者、木匠”。\textit{tekton} 源自原始印欧语 *\textit{teks-}“\textbf{编织、制造}”，其后代还包括 \textbf{text}（编织物 $\to$ 文本）、\textbf{texture}（纹理）、\textbf{technique/technology}（制作的技艺）。建筑师与纺织工、程序员共享一个“编织结构”的祖先词根。
\end{origin}

\vw{arrange}{/əˈreɪndʒ/}{vt. \& vi.}

\sense{vt. \& vi.}{整理，布置}
\ex{She arranged the flowers in a blue vase.}
\ex{The books are arranged alphabetically by author.}

\sense{vt.}{安排，筹备}
\col{arrange sth. for sb.}
\ex{The secretary arranged a meeting for the two managers.}
\col{arrange for sb. to do sth.（注意 for 的存在）}
\ex{We arranged for a car to pick you up at the airport.}
\col{arrange to do sth.}
\ex{We've arranged to meet at six outside the cinema.}

\begin{gram}[arrange that\dots 与“隐式虚拟”]
\textit{We arranged that the papers (should) be sent by express mail.} 我们安排文件用快递寄出。——\textbf{arrange / suggest / demand / insist / order} 等表“指令、要求”的动词后的 that 从句中用 \textit{(should+) 动词原形}；该 should 常可省略，从句表面看不出任何“虚拟”标记，故称隐式虚拟。名词形式 \textbf{arrangement}：\textit{make arrangements for} 为……做好安排。
\end{gram}

\begin{origin}
古法语 \textit{arengier}“排成一行”：\textit{a-} + \textit{rengier}“入列”，名词 \textit{reng}“行列”借自日耳曼语（与英语 \textbf{ring}、\textbf{rank} 同源，原始日耳曼语 *\textit{hringaz}“环、圈”）。士兵围成环站队 $\to$ 排列 $\to$ 安排。英语中一批带 a- 的法语借词（\textbf{assign, attire, amend}）都经历了同样的“动词加缀”改造。
\end{origin}

\vw{aside}{/əˈsaɪd/}{adv.}

\sense{adv.}{到旁边，在旁边；留，存}
\col{step / move / stand aside}
\ex{He stepped aside to let the old lady pass.}
\col{brush / push sth. aside（引申：不理会）}
\ex{He brushed aside all objections and went ahead with the plan.}
\col{joking / putting aside（说正经的，插入语）}
\ex{Joking aside, we really need to work harder.}
\col{set sth. aside（把……放一边；留出，拨出）}
\ex{Try to set aside half an hour every day for reading.}
\ex{A sum of money was set aside for the education of refugees.}
\col{aside from（= apart from）}
\ex{Aside from a few minor cuts, no one was hurt in the accident.}

\vw{association}{/əˌsəʊsiˈeɪʃn/}{n.}

\sense{n.}{联想，联系}
\col{the association of A with B}
\ex{The smell of chalk always carries associations of my school days.}
\col{in association with}
\ex{The concert was organized in association with the local radio station.}

\sense{n.}{协会，社团}
\ex{He is a member of the Photographers' Association.}

\sense{n.}{联合，交往}
\ex{Their close association with the firm lasted twenty years.}

\begin{origin}
拉丁语 \textit{socius}“\textbf{伙伴、盟友}”（原始印欧语 *\textit{sokw-yo-}，出自 *\textit{sekw-}“跟随”——伙伴即“跟随你的人”）。同根词 \textbf{social}、\textbf{society}、\textbf{associate}、\textbf{sociology}。与“跟随”同源的还有 \textbf{sequel}（续集）、\textbf{consequence}（随之而来者 $\to$ 后果）、\textbf{suit}（跟随者 $\to$ 随从$\to$诉讼；套装）。“社会”的本义就是“追随者们组成的群体”。
\end{origin}

\vw{atmosphere}{/ˈætməsfɪər/}{n.}

\sense{n.}{气氛，氛围}
\col{create / spoil a\dots atmosphere}
\ex{The restaurant has a cozy, relaxed atmosphere.}
\ex{Use candles to create a romantic atmosphere at dinner.}

\sense{n.}{大气（层）}
\ex{The rocket pierced the atmosphere and entered outer space.}

\sense{n.}{气体；空气}
\ex{The atmosphere in the mine was tested for poisonous gases.}

\begin{origin}
现代拉丁语科学造词 \textit{atmosphaera}，由希腊语 \textit{atmos}“\textbf{水汽、呼出的气}”（与动词 \textit{atmein}“呼气”同族）+ \textit{sphaira}“球”（\textbf{sphere}）拼合而成，17 世纪由炼金术士（一说开普勒）所造，字面义“包裹地球的\textbf{蒸汽球}”。它是“希腊零件 + 现代组装”的科学词典型，同类还有 \textbf{photosynthesis}、\textbf{telephone}。保留了希腊词尾 -os 的还有 \textbf{atmopheric} 的拼写。
\end{origin}

%================================
\section{B}
%================================

\vw{barren}{/ˈbærən/}{adj.}

\sense{adj.}{贫瘠的，不毛的}
\col{barren land / soil / hillsides}
\ex{Nothing grows on the barren slopes of the mountain.}

\sense{adj.}{不结果实的}
\ex{The old pear tree has been barren for years.}
\col{a barren discussion（引申：毫无成果的，书面语）}
\ex{The talks proved barren, and both sides went home angry.}

\begin{origin}
经由古法语 \textit{barhain}“不育的”进入英语，词源不明，一般认为是\textbf{前拉丁时期的地中海底层语言（substratum）词汇}——这与它 stubborn 的拼法（-rr-）相吻合。法语中同源词 \textit{barren} 在方言里仍指“不生育的母羊”。英语中来源成谜的日常词还有 \textbf{bad、big、dog}，它们常比“有家谱”的拉丁词更古老。
\end{origin}

\vw{bathe}{/beɪð/}{vt.}

\sense{vt.}{以（光线）洒满，覆盖，使沐浴（在光线里）}
\col{be bathed in light / sunshine / sweat（高频被动）}
\ex{At dawn the whole valley was bathed in golden light.}
\ex{His face was bathed in sweat after the race.}

\sense{vt.}{用水清洗（伤口等）}
\col{bathe a wound}
\ex{The nurse gently bathed the cut and applied a dressing.}

\begin{gram}[bathe / bath / breathe 的“清浊对应”]
英式英语中 \textbf{bathe} 作不及物动词可指“游泳”（旧式：\textit{go for a bathe}）；\textbf{bath} $n.$ 浴缸；洗澡。注意一组“名清动浊”的对子：\textbf{breath} /breθ/（名）vs.\ \textbf{breathe} /briːð/（动）、\textbf{bath} /bɑːθ/ vs.\ \textbf{bathe} /beɪð/——名词用清音 /θ/ 且元音较短，动词用浊音 /ð/ 且（往往）元音较长。这条规律覆盖 \textit{mouth, use, house} 等一批词，是听拼结合的记忆利器。
\end{gram}

\begin{origin}
古英语 \textit{bæðian}“浸洗”，与名词 \textit{bæth}（bath）同源，均可追溯至原始印欧语 *\textit{bhē-}“\textbf{温暖}”——“澡”的本义是“用温水暖身”。它的远亲是拉丁语 \textit{fovere}“温暖、孵化”（\textbf{foment} 热敷、煽动）与希腊语 \textit{tēkein}“融化”（\textbf{thaw} 融化与之平行）。“沐浴在阳光里”的意象，从词源起就是“被暖意包裹”。
\end{origin}

\vp{be down to sb./sth.}

\sense{v. 短语}{由……引起（或造成）（多用于口语归因）}
\ex{``Why is the shop closed?'' ``It's all down to the pandemic.''}
\ex{His success is largely down to sheer hard work.}
\col{It's down to sb. to do sth.（做某事是某人的责任）}
\ex{It's down to you to check the tickets.}
\col{from A down to B（下至、直到）}
\ex{Everyone, from the CEO down to the cleaner, attended the ceremony.}

\vw{belief}{/bɪˈliːf/}{n.}

\sense{n.}{看法，信念}
\col{a belief that\dots（同位语从句高频）}
\ex{She expressed her firm belief that things would improve.}
\col{It is my belief that\dots}
\ex{It is my belief that honesty is the best policy.}

\sense{n.}{信仰}
\col{religious / political beliefs}
\ex{People of all beliefs are welcome at the conference.}

\sense{n.}{相信，信心}
\col{belief in sb./sth.（介词固定用 in）}
\ex{She has lost all belief in her own abilities.}
\col{beyond belief（难以置信地）}
\ex{The mess they left behind was beyond belief.}
\col{disbelief（不信，怀疑）}
\ex{She stared at the letter in disb.elief.}

\begin{origin}
中古英语 \textit{beleave}，来自动词 \textit{belyven}：\textit{be-} + \textit{leven}“爱、视为可爱”（古英语 \textit{lēof}“亲爱的”，与德语 \textit{lieb} 同源，原始印欧语 *\textit{leubh-}“\textbf{爱、渴望}”）。“相信”的本义竟是“\textbf{在心里珍爱、看重某人}”。同根词 \textbf{love}、\textbf{lief}（乐意的，古雅）、德语 \textbf{Liebe}（爱）。belief 词尾的 -f 与 behave/half 一样，是 15 世纪为与复数动词形式（believe, halves）区分而人为“仿古”加上的，并无古老依据。
\end{origin}

\vw{bend}{/bend/}{vi. \& vt.\ (bent, bent)\ \&\ n.}

\sense{vi. \& vt.}{（使）拐弯，弯曲}
\ex{The road bends sharply to the right just after the bridge.}
\ex{He bent the wire into a circle.}
\col{bend down / over}
\ex{She bent down to pick up the coin.}
\col{bend the rules（通融、变通规则）}
\ex{The teacher sometimes bends the rules for hard-working students.}
\col{be bent on (doing) sth.（一心想要）}
\ex{He is bent on winning the championship at any cost.}

\sense{n.}{拐弯，弯道}
\ex{Slow down when you approach a blind bend.}

\begin{origin}
古英语 \textit{bendan}“（给弓上弦而）使其弯曲；用绳索捆缚”，名词 \textit{bend} 本指“\textbf{捆扎用的皮绳}”，源自原始印欧语 *\textit{bhendh-}“\textbf{捆绑}”。这个词根在日耳曼语中演化为 \textbf{bind}（捆绑）、\textbf{band}（带子；乐队“绑在一起的乐手”）、\textbf{bond}（纽带）、\textbf{bundle}（一捆）；在拉丁语中借入为 \textbf{bandage}（绷带）、\textbf{bend} 的使役义。“弯曲”与“捆缚”共用词根，因为捆缚即使之弯绕。
\end{origin}

\vw{bet}{/bet/}{n.\ \&\ vi.\&vt.\ (bet, bet)}

\sense{n.}{打赌，赌注；预计，估计}
\col{make / place a bet (on sth.)}
\ex{We made a bet on who would win the World Cup.}
\col{a good / safe bet（有把握的选择）}
\ex{Packing an umbrella is always a safe bet in London.}

\sense{vi. \& vt.}{下赌注，打赌}
\col{bet (sb.) (sth.) on\dots}
\ex{I bet him ten yuan that our team would lose.}

\sense{vt.}{敢说（口语）}
\col{I bet (that)\dots / You bet.}
\ex{I bet she's forgotten the meeting entirely.}
\ex{``Are you coming to the party?'' ``You bet!''}
\col{don't bet on it（别指望）}
\ex{``He'll apologize, won't he?'' ``Don't bet on it.''}

\begin{origin}
17 世纪突然出现的俚语词，来历成谜：一说来自 \textbf{abet}（教唆，古法语 \textit{abeter}“a- + 咬叫着驱使猎犬”），一说为 \textbf{better}“使更好”的赌徒变形——押注是为了“改善”赔率。它与 \textbf{beat}（打）形似却无关。词汇史中，赌博行话（\textbf{stake} 赌注、\textbf{odds} 赔率、\textbf{jackpot} 头奖）多为这种“来历不明的黑话”，因为俚语刻意逃避书面记录。
\end{origin}

\vw{bitter}{/ˈbɪtər/}{adj.}

\sense{adj.}{味苦的}
\ex{The medicine tasted so bitter that the child spat it out.}

\sense{adj.}{激烈的（争论、斗争）}
\col{a bitter dispute / quarrel / fight}
\ex{The two companies are locked in a bitter legal battle.}

\sense{adj.}{令人难过的}
\col{a bitter disappointment / blow / to the bitter end}
\ex{Losing the final was a bitter disappointment for the whole team.}
\ex{The defenders fought on to the bitter end.}

\sense{adj.}{严寒的}
\col{a bitter wind / night}
\ex{A bitter wind blew from the north that night.}

\begin{origin}
古英语 \textit{biter}，字面义“\textbf{咬人的}”，由 *\textit{bit-}（bite 的词干，原始印欧语 *\textit{bheid-}“劈裂”）+ 工具后缀 -er 构成——苦味正是“\textbf{咬}舌头”的味道，与汉语“苦得咬嘴”异曲同工。同根词 \textbf{bite}、\textbf{bit}（马嚼子；一点“咬下来的一口”）、\textbf{bittern}（麻鸦，苦味鸟？实为拟声）。啤酒术语 \textbf{bitters}（苦味酒）与 \textbf{bitter end}（“缆绳末端”，航海转义）都是它的后代。
\end{origin}

\begin{gram}[味觉通感]
sweet（甜$\to$悦人的、温柔的）、\textbf{bitter}（苦$\to$痛苦的、激烈的、严寒的）、sour（酸$\to$乖戾的，\textit{go sour} 变糟）、salty（咸$\to$〔俚〕风趣的）、hot（辣$\to$辛辣的、热门的）。这种由味觉向情绪、温度、程度引申的现象是英语形容词多义的主要来源之一，阅读时应先判断修饰对象的类型再取义。
\end{gram}

\vw{blame}{/bleɪm/}{vt.\ \&\ n.}

\sense{vt.}{把……归咎于，责怪，指责}
\col{blame sb. for sth.}
\ex{Don't blame yourself for the accident --- it wasn't your fault.}
\col{blame sth. on sb./sth.}
\ex{The railway company blamed the delay on the heavy snow.}
\col{be to blame (for sth.)}
\ex{Nobody was to blame for the breakdown of the talks.}

\sense{n.}{责任，责备，指责}
\col{put / lay the blame (for sth.) on sb.}
\ex{The government tried to put the blame for inflation on world markets.}
\col{take / carry / accept the blame for}
\ex{The captain took theblame for the team's defeat.}

\begin{gram}[be to blame：主动形式表被动意义]
“某人该受责备”须说 \textit{sb. is to blame}，而非\textit{sb. is to be blamed}（惯用法上被视为错误）。这是“be + 不定式”结构的特殊用法：不定式主动式作表语表达“应尽的责任或安排”，此处语义上“某人应受责备”却用主动不定式，属于英语中少数\textbf{主动表被动}的情形（类似的还有 \textit{The house is to let.} 此房出租）。考试中该结构几乎逢考必有。
\end{gram}

\begin{origin}
经由古法语 \textit{blasmer}“指责”进入英语，而 \textit{blasmer} 是晚期拉丁语 \textit{blasphemare}“\textbf{亵渎、出言不敬}”的讹变——没错，blame 与 \textbf{blasphemy}（亵渎神明）是同一个词，都出自希腊语 \textit{blasphēmein}：\textit{blasphēmos}“恶语的”+ \textit{phēmē}“言语”（原始印欧语根 *\textit{bhā-}“说话”，同根 \textbf{fame} 名声、\textbf{fate} 神谕“被说出的定数”、\textbf{infant} 不会说话的人）。“责备”在词源上是一次“小小的渎神”。
\end{origin}

\vw{blues}{/bluːz/}{n.}

\sense{n.}{the blues 布鲁斯音乐，蓝调；忧郁，悲伤}
\ex{Billie Holiday is one of the greatest blues singers of all time.}
\col{feel blue / have the blues（感到忧郁，口语）}
\ex{He's been having the blues since his best friend moved away.}

\begin{origin}
音乐义来自 19 世纪美国黑人短 语 \textit{the blue devils}“\textbf{蓝色魔鬼}”（幻觉与抑郁，其单数 \textbf{blue devil} 在 17 世纪指酒后的极度沮丧）。“蓝色 = 忧郁”是英语独有的色彩隐喻（对比汉语“红”之兴旺），同族隐喻构成一个系统：\textbf{green with envy}（嫉妒得发绿）、\textbf{see red}（大怒）、\textbf{in the pink}（非常健康）、\textbf{black mood}（阴郁情绪）。语言学的“萨丕尔–沃尔夫假说”常以此为喻：色彩的情感义是文化赋予的，并非普遍真理。
\end{origin}

\vw{bond}{/bɒnd/}{n.}

\sense{n.}{纽带，联系}
\col{a bond between A and B / a bond with sb.}
\ex{Years of working together forged a strong bond between the two families.}
\col{strengthen / break the bonds}
\ex{Nothing could break the bond between the mother and her children.}

\sense{n.}{连接，结合}
\ex{In chemistry, two hydrogen atoms and one oxygen atom form water through chemical bonds.}
\col{bond with sb.（动词：与某人建立信赖关系）}
\ex{Fathers bond with their babies through play.}
\col{government bonds（债券）}
\ex{Grandpa keeps part of his savings in government bonds.}

\begin{origin}
\textbf{bond} 是 \textbf{band}“带子”的变体（见 bend 词条），四词 \textit{bind--band--bond--bundle} 是原始印欧语 *\textit{bhendh-}“捆绑”在英语里的“一鱼四吃”。债券称 bond，因为它是“\textbf{捆}住双方义务的法律文书”；化学键是 20 世纪对 bond“连接”义的科学借用；007 系列 \textit{Bond} 与此无关（本为鸟名）。
\end{origin}

\vw{boom}{/buːm/}{n.\ \&\ vi.}

\sense{n.}{（贸易和经济活动的）激增，繁荣}
\col{an economic / tourist / baby boom}
\ex{The tourism boom brought jobs and investment to the coastal town.}

\sense{vi.}{迅速发展，激增，繁荣昌盛}
\ex{Business is booming in the southern manufacturing cities.}
\col{boom and bust（繁荣与萧条的交替）}
\ex{The region's economy has gone through repeated cycles of boom and bust.}
\col{boom（拟声：低沉轰鸣）；a booming voice（洪亮的嗓音）}
\ex{Thunder boomed in the distance.}

\begin{origin}
拟声词，源自荷兰语水手用语 \textit{boom}“\textbf{帆桁、横杆}”（因撞击船身发出“砰”声），本义与英语 \textbf{beam} 同源（原始日耳曼语 *\textit{baumaz}“树、梁”，荷兰语 \textit{boom} 至今兼指“树”）。“经济激增”义出自 19 世纪美国报纸：\textit{boom} 描述枪声、股价“轰然窜升”的声音，遂成名 词。防波的“拦河木” \textbf{boom}（log boom）保留“横木”本义。
\end{origin}

\vp{break with sth.}

\sense{v. 短语}{与某事终止关联，破除（传统、惯例）}
\ex{The new manager broke with tradition and promoted staff purely on merit.}
\ex{The country broke with its isolationist past after the war.}
\col{break with sb.（与某人决裂）}
\ex{He broke with his business partner over the new project.}

\vw{breath}{/breθ/}{n.}

\sense{n.}{呼吸的空气；一次吸入的空气}
\col{take a deep breath / catch one's breath}
\ex{Take a deep breath and jump into the water.}
\ex{She stopped to catch her breath at the top of the stairs.}
\col{out of breath / hold one's breath}
\ex{He was out of breath after chasing the bus.}
\ex{The whole stadium held its breath as the ball flew toward the goal.}
\col{take sb.'s breath away（令人惊叹）}
\ex{The view from the summit took our breath away.}

\sense{n.}{微量，迹象}
\col{not a breath of wind / air}
\ex{There wasn't a breath of wind on the lake that morning.}

\begin{gram}[breath / breathe：名清动浊]
\textbf{breath} /breθ/ 名词 vs.\ \textbf{breathe} /briːð/ 动词：名词短元音 + 清音 /θ/，动词长元音 + 浊音 /ð/（详见 bathe 词条）。同类对：\textit{bath/bathe}、\textit{cloth/clothe}、\textit{mouth}(n.)/(v.)。听写与听力辨音的高频陷阱。
\end{gram}

\begin{origin}
古英语 \textit{brǣth}“呼出的气、气味”，与动词 \textit{brǣthan}（breathe）同源，出自原始印欧语 *\textit{breh\u{1}-}“\textbf{散发气味、蒸发}”，其近亲是拉丁语 \textit{fragrare}“飘香”（\textbf{fragrant} 芬芳的）。“呼吸”在词源上竟是“\textbf{发气味}”——古人认为灵魂随气息出入，故 \textbf{s[l.irit}（拉丁语 \textit{s[l.iritus}“一口气”）与 \textbf{ins[l.ire}（吹入气息$\to$赋予灵感）、\textbf{expire}（气息离体$\to$断气；期满）同属“气息 = 生命”的隐喻家族。
\end{origin}

\vw{breeze}{/briːz/}{n.}

\sense{n.}{微风，和风}
\ex{A gentle breeze stirred the curtains at the open window.}
\col{a breeze（口语：轻而易举的事）}
\ex{The exam was a breeze.}
\col{breeze through sth.（轻松完成）}
\ex{She breezed through the interview and got the job.}

\begin{origin}
16 世纪出现，词源成谜，多认为借自西班牙比斯开湾水手的 \textbf{brisal/briza} 类航海词；西班牙语 \textit{brisa} 至今指海风。海员用语中， breeze 特指\textbf{白天从海面吹向陆地}的向岸风（与夜间的 land breeze 相对），“微风”义是陆地化后的弱化。俚语义“小菜一碟”出现在 20 世纪初的美国校园。
\end{origin}

\vw{brilliant}{/ˈbrɪliənt/}{adj.}

\sense{adj.}{聪颖的，技艺高的}
\ex{Einstein was one of the most brilliant minds of the twentieth century.}

\sense{adj.}{巧妙的}
\ex{That's a brilliant idea --- why didn't I think of it?}

\sense{adj.}{很成功的（英式口语）}
\ex{``I've got us two tickets for the final!'' ``Brilliant!''}

\sense{adj.}{明亮的}
\ex{The lake glittered under the brilliant afternoon sun.}
\col{brilliantly（副词）}
\ex{The plan worked brilliantly.}

\begin{origin}
法语 \textit{brillant}“闪耀的”，动词 \textit{briller} 来自意大利语 \textit{brillare}，通说源自拉丁语 \textit{beryllus}“\textbf{绿柱石、海蓝宝石}”（又借自希腊语 \textit{bēryllos}，本或指印度某产地）。“宝石般闪光”是其本义，故同根词 \textbf{beryl}（绿柱石）与 \textbf{brilliant} 曾指同一种晶亮的石头；“brilliant cut”（明亮式切工）至今是钻石琢型的术语，而“ brilliants”在旧小说里可直接指钻石。
\end{origin}

%================================
\section{C}
%================================

\vw{calling}{/ˈkɔːlɪŋ/}{n.}

\sense{n.}{职业，事业（尤指视为使命的工作）}
\ex{Teaching is more than a job --- for her it is a true calling.}

\sense{n.}{使命感}
\col{have a calling for / to do sth.}
\ex{He felt a calling to serve his community as a doctor.}

\begin{lex}[job / career / calling / vocation]
job（谋生的工作）$<$ career（长期事业）$<$ calling（奉为使命的天职）。\textbf{vocation} 与 calling 近义，词根 \textit{voc-} 表“召唤”（见 advocate 词条）。“职业”词义由“被（上天）呼唤”的宗教隐喻而来，这一隐喻在汉语“天职”“神职”中同样存在。
\end{lex}

\vw{calorie}{/ˈkæləri/}{n.}

\sense{n.}{卡，卡路里；千卡（食物热量单位）}
\col{high-calorie / low-calorie / calorie intake}
\ex{A big plate of fried rice contains about 600 calories.}
\ex{You need to burn more calories than you take in to lose weight.}

\begin{origin}
19 世纪法国科学家为热学新造的单位词，直接取自拉丁语 \textit{calor}“\textbf{热量}”（原始印欧语 *\textit{kelə-}“温暖”，同根 \textbf{cauldron} 大锅“烧热之器”、\textbf{luke-warm} 的 \textit{lew}“微温”、\textbf{nonchalant}“不在乎”$\leftarrow$“不发热的”）。有意思的是，\textbf{calculate}（计算）与之无关：它来自 \textit{calculus}“\textbf{小石子}”（古人以石子计数），所以“算卡路里”是拿两颗不同的“石头”在算。
\end{origin}

\vw{campaign}{/kæmˈpeɪn/}{n.\ \&\ vt.\&vi.}

\sense{n.}{（宣传）活动，运动}
\col{launch / wage a campaign for / against sth.}
\ex{The city has launched a campaign against food waste in schools.}
\col{an advertising / election campaign}
\ex{Her election campaign focused on health care and education.}

\sense{n.}{战役}
\ex{The general led a brilliant defensive campaign in the mountains.}

\sense{vt. \& vi.}{发起运动，参加活动}
\col{campaign for / against sth.}
\ex{Local residents are campaigning against the construction of the new highway.}
\ex{She s[l.ent years campaigning for women's right to vote.}

\begin{origin}
法语 \textit{campagne}“\textbf{开阔的乡野}”，借自意大利语 \textit{campagna}“平原”，最终来自拉丁语 \textit{campus}“\textbf{田野、平原}”（同根 \textbf{camp} 营地、\textbf{champion} 原指“战场上的斗士”、\textbf{Cambridge}“格兰河畔的营地”之一解）。军事上“把军队\textbf{拉到野外作战}的一个季节”称为 campaign，故 17 世纪英语 campaign 本指“战役”；“宣传运动”是 19 世纪借“一系列攻势”的比喻。
\end{origin}

\vw{career}{/kəˈrɪər/}{n.}

\sense{n.}{经历，事业；生涯，职业}
\col{a career in politics / medicine / journalism}
\ex{She started her career as a junior reporter.}
\col{career prospects / a career move}
\ex{Moving to Shanghai was a smart career move.}
\col{at the peak / height of one's career}
\ex{At the height of his career, he suddenly announced his retirement.}

\begin{origin}
法语 \textit{carrière}“\textbf{赛车道}”，借自意大利语 \textit{carriera}“车行道”，源自拉丁语 \textit{carraria (via)}“运货车的路”，来自 \textit{carrus}“四轮马车”（凯尔特语借词，同根 \textbf{car}、\textbf{carry}、\textbf{charge}“装上马车$\to$装载；冲锋；收费”、\textbf{cargo}）。16 世纪英语 \textit{career} 就指“飞驰；马术竞赛跑道”，短语 \textit{in full career}“全速飞奔”保留此义；“人生赛道”的比喻在 19 世纪定型——\textbf{职业生涯}本就是一场 \textit{car race}。注意：另一个同音词 \textit{korea} 拼法不同。
\end{origin}

\vw{cast}{/kɑːst/}{vt.\ (cast, cast)\ \&\ n.}

\sense{vt.}{投射（光线、阴影）；向……投以（视线、笑容等）}
\col{cast light / a shadow on sth.（本义与喻义）}
\ex{The setting sun cast long shadows across the fields.}
\ex{The scandal cast a shadow over his final year in office.}
\col{cast a glance / an eye over sth.}
\ex{She cast a quick glance over her shoulder.}

\sense{vt.}{投，抛}
\ex{The fishermen cast their nets into the sea at dawn.}

\sense{vt.}{选派角色}
\ex{The director cast an unknown actress in the leading role.}

\sense{n.}{全体演员；投，抛}
\ex{The whole cast deserve praise for this wonderful production.}
\col{cast your mind back to sth.（回想起）}
\ex{Cast your mind back to your first day at school.}

\begin{gram}[三者同形的不规则动词]
cast--cast--cast（同型者：cut, hit, hurt, let, put, set, shut, split, cost）。这类词的过去分词与原形同形，因此在现在完成时中易与被动语态混淆：\textit{The net was cast.}（被动，网被撒下）vs.\ \textit{He has cast the net.}（完成时）。判断依据：句中有无 by 逻辑主语、主语是施动还是受动。
\end{gram}

\begin{origin}
借自古诺斯语（维京人的北欧语）\textit{kasta}“投掷”，13 世纪进入英语，属\textbf{北欧借词}层（与 sky, skin, they, take 同批，丹麦法区 Danelaw 的遗产）。“选角”义来自“把演员\textbf{掷入}角色”的戏剧行话；“铸型”义（\textit{cast iron} 铸铁）来自“把熔化的金属\textbf{浇投}进模具”；\textbf{forecast}（fore + cast“往前掷”$\to$预报）则是 18 世纪的新组合。
\end{origin}

\vw{characterize}{/ˈkærəktəraɪz/}{vt.}

\sense{vt.}{是……的特征，以……为典型}
\col{be characterized by sth.（被动式最常见）}
\ex{The early poems are characterized by simplicity and quiet sadness.}

\sense{vt.}{使……具有特点}
\ex{A s[l.irit of innovation characterizes the company's culture.}

\sense{vt.}{描述，刻画}
\ex{The novel vividly characterizes life in a small fishing village.}

\begin{origin}
希腊语 \textit{kharaktēr}“\textbf{刻上去的印记}”，来自动词 \textit{kharassein}“刻、削尖”（名词 \textit{kharax}“尖头木桩”）。古希腊人刻章记功、铸币压模，“\textbf{压印出的标记}”遂成一个人区别于他人的“印记”，即性格；在印刷术中又成“\textbf{铸出的符号}”，即字符。一个词根串起“性格”“角色”“字符”“刻画”四义。同根 \textbf{character}、\textbf{characteristic}。
\end{origin}

\vw{chorus}{/ˈkɔːrəs/}{n.}

\sense{n.}{合唱曲；副歌}
\ex{The audience joined in the chorus of the familiar song.}

\sense{n.}{合唱团}
\ex{She sings in the school chorus.}
\col{in chorus（齐声地）}
\ex{``Good morning, teacher!'' the students replied in chorus.}
\col{a chorus of praise / protest（众口一词的赞扬／抗议）}
\ex{The decision was greeted with a chorus of disapproval.}

\begin{origin}
希腊语 \textit{choros}“\textbf{环舞、歌舞队}”（原始印欧语 *\textit{gher-}“抓、围圈”，同根 \textbf{garden}“围起来的园子”、\textbf{yard} 庭院）。古希腊戏剧的歌队围成圆场且舞且歌，故“合唱”与“圆圈”同源。拉丁语转写为 \textit{chorus}，其“教堂唱诗班”义生出英语 \textbf{choir}（历史上拼作 \textit{quire}，与 \textit{chor-} 为同一词的不同“外语口音”版本）；圣诞颂歌 \textbf{carol} 最初也是“环舞之歌”。
\end{origin}

\vw{claim}{/kleɪm/}{n.\ \&\ vt.}

\sense{n.}{声明，宣称；所有权；索赔}
\col{make a claim / make a claim on/against sth.}
\ex{He made wild claims about his own importance that nobody believed.}
\ex{You can make a claim on your insurance for the damaged goods.}
\col{lay claim to sth.（声称对……的所有权）}
\ex{Three museums lay claim to the same painting.}

\sense{vt.}{宣称，声称}
\col{claim that\dots / claim to do / to have done sth.}
\ex{He claims to have met the president, but I doubt it.}

\sense{vt.}{要求；索取}
\ex{Passengers should claim their luggage at Carousel 3.}

\sense{vt.}{获得；夺去（生命）（教材未列，极高频）}
\col{claim sb.'s life / claim lives}
\ex{The earthquake claimed over two hundred lives.}

\begin{gram}[claim to have done：不定式完成式]
\textit{claim / seem / appear / pretend + to have done} 表“声称／似乎／假装\textbf{已经}做过”，完成式体现动作先于谓语发生。与从句转换：\textit{He claims to have met her.} $=$ \textit{He claims that he \textbf{has met} her.} 注意 \textit{He pretended to be reading.}（进行式，假装“正在”）与完成式的体貌差别——不定式的“时”信息全部由 have / be 体现。
\end{gram}

\begin{origin}
拉丁语 \textit{clamare}“\textbf{大声呼喊}”（原始印欧语 *\textit{kel-}“喊叫”，其扩展形式 *\textit{kal-} 给出 \textbf{call}“呼叫”——两者是同一词根的拉丁、日耳曼两支后裔）。\textit{claim} 经由古法语 \textit{clamer} 进入英语，本为法律用语“高声宣布权利”。同根词 \textbf{clamor}（喧嚷）、\textbf{exclaim}（喊出）、\textbf{proclaim}（公开宣布）、\textbf{acclaim}（欢呼喝彩）、\textbf{reclaim}（收回）。既然“主张权利”全靠“喊”，难怪 claim 同时是“宣称”与“认领”。
\end{origin}

\vw{clothing}{/ˈkləʊðɪŋ/}{n.}

\sense{n.}{衣服，服装（总称）}
\col{an item / a piece of clothing（不可数，借 item/piece 计数）}
\ex{Warm clothing is essential for the expedition to Antarctica.}
\col{protective / waterproof / second-hand clothing}
\ex{The Red Cross distributed clothing to the flood victims.}

\begin{lex}[clothing / clothes / cloth]
\textbf{clothing} 不可数，是“衣物”的正式总称，涵盖衣帽鞋袜；\textbf{clothes} 只有复数形式（无单数！），谓语用复数：\textit{These clothes are dirty.}，计数用 \textit{an article of clothes}；\textbf{cloth} 不可数，指“布料”（\textit{a cloth} 则指一块抹布）。三词同源异用，是名词可数性的经典考点。
\end{lex}

\vw{clue}{/kluː/}{n.}

\sense{n.}{线索，提示}
\col{a clue to / as to sth.（固定介词 to）}
\ex{The police found an important clue to the identity of the thief.}
\col{not have a clue（毫无头绪，口语）}
\ex{``Where are the keys?'' ``I haven't a clue.''}

\sense{n.}{迹象}
\ex{The dusty bookshelf was a clue that nobody had lived there for years.}

\begin{origin}
\textbf{clue} 是 \textbf{clew}“线团”的变体拼写，出自希腊神话\textbf{忒修斯与迷宫}：克里特公主阿里阿德涅给忒修斯一团线（\textit{clew of thread}），让他杀牛头怪后循线走出迷宫。16 世纪起“\textbf{引你走出迷宫的那根线}”就成了“线索”。莎士比亚《特洛伊罗斯与克瑞西达》仍拼作 \textit{clew}。一词的词源即一个完整的故事——这是英语“神话遗产词”的经典代表。
\end{origin}

\vw{combination}{/ˌkɒmbɪˈneɪʃn/}{n. \& combine /kəmˈbaɪn/ vt.\&vi.}

\sense{n.}{结合体，联合体}
\col{a combination of A and B / a combination of factors}
\ex{Success results from a rare combination of talent and hard work.}
\ex{The failure was due to a combination of bad luck and poor judgment.}

\sense{n.}{联合，混合}
\ex{The color purple is a combination of red and blue.}

\sense{combine vt.\&vi.}{（使）融合；（使）结合，（使）混合；兼做；合并}
\col{combine A with/and B}
\ex{The film combines humor with a serious message about the environment.}
\col{combine business with pleasure（工作娱乐两不误，习语）}
\ex{The conference in Hawaii combined business with pleasure.}

\begin{origin}
晚期拉丁语 \textit{combinare}“两两相合”：\textit{com-}（一起）+ \textit{bini}“\textbf{成对的}”，\textit{bini} 是 \textit{bis}“两次”的形容词形（原始印欧语 *\textit{dwo-}“二”，同根 \textbf{two}、\textbf{binary} 二进制、\textbf{bicycle} 双轮车）。所以 combine 的字面义是“\textbf{把二放在一起}”。令人意外的亲戚：\textbf{biscuit}（bis + coctus“烤两次”$\to$烤两遍的饼）、\textbf{bini} 还潜藏在 \textbf{combination lock}“两两对号的密码锁”里。
\end{origin}

\vw{complex}{/ˈkɒmpleks/}{adj.\ \&\ n.}

\sense{adj.}{复杂的，难懂的}
\col{a complex problem / issue / system}
\ex{The human brain is far more complex than any computer.}

\sense{n.}{建筑群；复合体}
\ex{A new shopping and leisure complex will open near the station.}

\sense{n.}{情结（心理学术语，教材未列，常见于阅读）}
\col{an inferiority / superiority complex}
\ex{His fear of public speaking stems from an inferiority complex.}

\begin{lex}[complex / complicated]
\textbf{complex} 中性词，指“由多部分构成的、内在结构复杂”，常含“精密”义（complex system）；\textbf{complicated} 常带贬义，指“令人费解的、不必要的繁琐”（complicated rules, a complicated relationship）。科学语篇用 complex，抱怨吐槽用 complicated。名词 \textbf{complexity}（复杂性）。
\end{lex}

\begin{origin}
拉丁语 \textit{complexus}“\textbf{编织在一起的、交织物}”，过去分词作名词，来自动词 \textit{complecti}：\textit{com-}（一起）+ \textit{plex-}（词根 \textit{plectere}“\textbf{编、辫}”，原始印欧语 *\textit{plek-}“编辫子”）。“复杂”的本义是“\textbf{多股线绞在一起}”。同根词 \textbf{plait}（辫子）、\textbf{ply}（纱层；股）、\textbf{perplex}（per- 彻底 + 编$\to$完全缠住$\to$使困惑）、\textbf{plexus}（神经丛）、\textbf{multiplex}（多股的、多厅影院）。荣格用 \textit{Komplex}“情结”取“情绪与观念\textbf{绞缠}成团”之意。
\end{origin}

\vw{compose}{/kəmˈpəʊz/}{vt.\ \&\ vi.}

\sense{vt. \& vi.}{作曲}
\ex{Mozart composed his first symphony at the age of eight.}

\sense{vt.}{组成（用于 be composed of）}
\col{be composed of\dots（= consist of，但可用被动）}
\ex{Water is composed of hydrogen and oxygen.}

\sense{vt.}{撰写}
\ex{She composed a short s[l.eech for the graduation ceremony.}

\sense{vt.}{使镇静}
\col{compose oneself}
\ex{He took a moment to compose himself before answering the difficult question.}
\col{composed / composure（镇静的／沉着）}
\ex{She remained composed throughout the interview.}
\ex{The captain kept his composure in the crisis.}

\begin{origin}
拉丁语 \textit{componere}“\textbf{放在一起}”：\textit{com-} + \textit{ponere}“放置”（原始印欧语 *\textit{po-}，同根的巨大家族：\textbf{position} 位置、\textbf{compound} 混合物、\textbf{postpone} 往后放$\to$推迟、\textbf{opponent} 置于对面者$\to$对手、\textbf{component} 放在一起的东西$\to$组件）。作曲、作文、使镇静共享“\textbf{把纷乱的元素安放归位}”这一意象；\textbf{compose oneself} 即“把心绪各就各位”。合成词 \textbf{composure} 字面义竟是“安放妥当的状态”。
\end{origin}

\vw{consist}{/kənˈsɪst/}{vi.}

\sense{vi.}{consist of\dots 由……组成（或构成）}
\ex{The committee consists of ten members from all walks of life.}
\ex{A typical English breakfast consists of eggs, bacon and toast.}
\col{consist in（在于，书面语）}
\ex{Happiness consists in contentment.}

\begin{gram}[状态动词 consist 的三条限制]
\textbf{consist of} 表静态的“构成”，因此：（1）不用进行时（\ding{55} \textit{is consisting of}）；（2）不用被动语态（\ding{55} \textit{is consisted of}，须改用 \textit{be composed of} 或 \textit{be made up of}）；（3）与 know, belong, contain 等状态动词一样描述恒常关系而非动作。改写练习：\textit{The team consists of 20 players.} $\Rightarrow$ \textit{The team \textbf{is made up of} / \textbf{is composed of} 20 players.}
\end{gram}

\begin{origin}
拉丁语 \textit{consistere}“\textbf{站在一起、站立不动}”：\textit{con-}（共同）+ \textit{sistere}“使站立”（\textit{stare}“站”的使役动词，原始印欧语 *\textit{stā-}“站立”——人类最稳定的姿态）。“由……组成”的字面义是“（与这些成员）\textbf{一同站稳}”。这个词根 \textit{stā-} 在英语中枝繁叶茂：\textbf{exist}（ex- 站出来$\to$显现$\to$存在）、\textbf{insist}（站在上面$\to$坚持）、\textbf{persist}（一直站着$\to$坚持不懈）、\textbf{resist}（站回来$\to$抵抗）、\textbf{stable}、\textbf{station}、\textbf{statue}。背下 consist，等于拿到半个家族的钥匙。
\end{origin}

\vw{constantly}{/ˈkɒnstəntli/}{adv.}

\sense{adv.}{始终，一直}
\ex{The refugee family lived in constant fear, constantly on the move.}

\begin{gram}[constantly + 进行时：情感色彩的“进行体”]
\textit{He is \textbf{constantly} leaving his dirty socks on the floor.} 他老是把脏袜子扔在地上。——进行时与 always, constantly, forever 连用（注意用进行时而非一般时），表达说话人对频繁动作的\textbf{厌烦或抱怨}：\ding{55} \textit{He constantly leaves\dots}（语义中性）vs.\ \checkmark\ \textit{He is constantly leaving\dots}（含抱怨）。
\end{gram}

\begin{origin}
\textbf{constant} 来自拉丁语 \textit{constans}“站得稳的”——\textit{con-} + \textit{stare}“站”（与 consist 同根，见上条）。“恒定”就是“\textbf{一直站着不倒}”。物理常数 \textit{constant}（重力常数）与化学的 \textbf{invariant} 都是“站稳”的后代。副词 -ly 后缀是古英语 \textit{-līċe}“以……的方式”的遗存，与形容词 -ly（\textit{-līċ}“具有……样子的”）本是同一个词的两种语法化。
\end{origin}

\vw{consume}{/kənˈsjuːm/}{vt.}

\sense{vt.}{吃，喝，饮（正式用语）}
\ex{The guests consumed vast quantities of cake and ice cream at the party.}

\sense{vt.}{消耗，耗费（尤指燃料、能量或时间）}
\col{consume energy / fuel / time / resources}
\ex{The new air conditioner consumes 30 percent less electricity than the old one.}
\ex{Writing the report consumed most of his weekend.}
\col{be consumed with jealousy / guilt（被……吞噬，文学用法）}
\ex{He was consumed with jealousy at his brother's success.}

\begin{gram}[sume 词族]
\textbf{consumer}（\textit{consumer rights} 消费者权益）；\textbf{consumption}（\textit{energy consumption} 能耗）。拉丁语 \textit{sumere}“取”的家族在英语中全员到齐：\textbf{assume}（取来据为己有$\to$假定；承担）、\textbf{presume}（预先取$\to$擅自认为）、\textbf{resume}（重新取回$\to$恢复；摘要）、\textbf{consume}（完全取走$\to$消耗）。反义场景词 \textbf{produce}（生产）$\leftrightarrow$ consume（消耗）。
\end{gram}

\begin{origin}
拉丁语 \textit{consumere}“\textbf{全部取走、耗尽}”：\textit{con-}（加强语气，完全）+ \textit{sumere}“拿取”，后者源自原始印欧语 *\textit{em-}“拿”（希腊语同根词 \textit{aneinai}“带走”，\textbf{example} 一词中 \textit{ex-}+\textit{emere}“取出者$\to$样品”亦出于此；\textbf{premium}“额外取走之物$\to$奖金、保费”）。罗马人的经济动词 \textit{consumere} 在 \textit{summa}“取的总数$\to$最高点、总额”里还藏着 \textbf{sum}（总和）与 \textbf{summit}（顶峰）。
\end{origin}

\vw{contradictory}{/ˌkɒntrəˈdɪktəri/}{adj.}

\sense{adj.}{相互矛盾的，对立的}
\col{contradictory statements / evidence / reports}
\ex{The two witnesses gave contradictory accounts of the accident.}

\begin{gram}[“自相矛盾”的英语表达链]
动词 \textbf{contradict}（与……矛盾：\textit{The two reports contradict each other.}；反驳某人：\textit{Don't contradict your mother!}）$\to$ 形容词 \textbf{contradictory} $\to$ 名词 \textbf{contradiction}（\textit{in contradiction to\dots} 与……相矛盾）。写作模板：\textit{Far from being contradictory, the two views are complementary.} 这两种观点非但不矛盾，反而互补。
\end{gram}

\begin{origin}
拉丁语 \textit{contradicere}“\textbf{逆着说}”：\textit{contra-}“相对、相反”（原始印欧语 *\textit{kom-tero-}“并置的”，与 \textbf{country}“对面之地”同根）+ \textit{dicere}“说”（*\textit{deik-}“指出”，同根 \textbf{dictate} 口述、\textbf{dictionary}“说话之书”$\to$词典、\textbf{predict} 预先说$\to$预言、\textbf{verdict} 如实地说$\to$裁决）。“矛盾”就是“\textbf{反着说话}”，与汉语“自相矛盾”（矛与盾两件兵器）取象完全不同，却同样生动。
\end{origin}

\vw{crisis}{/ˈkraɪsɪs/}{n.\ (pl.\ crises)}

\sense{n.}{危机，危急关头；危难时刻，病危期}
\col{an economic / financial / political crisis}
\ex{The country is facing its worst economic crisis in decades.}
\col{a crisis point / at the height of the crisis}
\ex{Talks reached crisis point when both sides walked out.}
\ex{The patient was in crisis for two days before the fever broke.}

\begin{gram}[希腊源 -is 名词的复数]
crisis $\to$ crises /ˈkraɪsiːz/，同族：\textbf{emphasis} $\to$ emphases、\textbf{analysis} $\to$ analyses、\textbf{basis} $\to$ bases、\textbf{hypothesis} $\to$ hypotheses、\textbf{diagnosis} $\to$ diagnoses、\textbf{oasis} $\to$ oases、\textbf{thesis} $\to$ theses。规律：-\textit{is} /ɪs/ $\to$ -\textit{es} /iːz/。更古老的第三变格还有 \textbf{criterion} $\to$ criteria、\textbf{phenomenon} $\to$ phenomena。发音陷阱：bases 既是 basis 的复数（/ˈbeɪsiːz/）又是 base 的复数（/ˈbeɪsɪz/），听 -s 的清浊可辨。
\end{gram}

\begin{origin}
希腊语 \textit{krisis}“\textbf{决断、病情转折点}”，与 \textbf{critic} 同根（\textit{krinein}“甄别、判断”，见下条）。希波克拉底用它指“疾病\textbf{见分晓}的日子”——转好或恶化的关口。英语借入后先作医学术语（17 世纪），再泛化为“紧要关头”。注意：汉语“\textbf{危}机”含“危险 + 机会”的民间拆字义，但希腊语原义只有“\textbf{判断、分晓}”，“机会”义是汉语的再创造。
\end{origin}

\vw{critic}{/ˈkrɪtɪk/}{n.}

\sense{n.}{批评家，评论家}
\ex{The play opened to mixed reviews from leading theater critics.}

\sense{n.}{批评者，挑剔的人}
\col{a critic of sb./sth.}
\ex{He has always been a harsh critic of online gaming.}
\col{be critical of（形容词：对……挑剔）}
\ex{Some employers are overly critical of new graduates.}
\col{criticize sb. for doing sth. / constructive criticism}
\ex{Teachers should criticize the work, not the person, and offer constructive criticism.}

\begin{origin}
希腊语 \textit{kritēs}“\textbf{裁判、判断者}”，来自 \textit{krinein}“\textbf{分开、甄别、决断}”（原始印欧语 *\textit{krei-}“筛分”）。这个词根是希腊留给现代学术的枢纽：\textbf{crisis}（决断时刻——见下条）、\textbf{criterion}（判断的依据$\to$标准）、\textbf{criticism}（甄别之学）、\textbf{endocrine}（内分 secret“筛进去”$\to$内分泌）、\textbf{recrimination}（反诉）。“批评家”在词源上不是“挑刺的人”，而是“\textbf{做甄别判断的人}”——批评本是判断力。
\end{origin}

\vw{crop}{/krɒp/}{n.}

\sense{n.}{庄稼，作物；收成，产量}
\col{grow / plant crops}
\ex{The main crops of the region are wheat, corn and soybeans.}
\col{a good / poor / record crop}
\ex{A record crop of apples was harvested this autumn.}
\col{cream of the crop（出类拔萃者）}
\ex{Only the cream of the crop of applicants gets into that university.}
\col{a crop of（一批新出现的人或物）}
\ex{A new crop of graduates enters the job market every July.}

\begin{origin}
古英语 \textit{cropp}“\textbf{植物的顶端、嫩芽}”（鸟的\textbf{嗉囊}亦谓之 crop——鼓起的顶袋）。“庄稼／收成”义由“\textbf{枝头之果}”引申而来，故动词 \textit{crop}“啃吃草尖”、\textit{cropped hair}“剪短的头发”皆取“去顶”之义；摄影的 \textbf{crop}（裁剪）同样来自“剪去顶梢”。短语 \textit{crop up}“意外出现”即“\textbf{从地里冒出头}”——庄稼的比喻至今活着。
\end{origin}

\vp{cry out for}

\sense{v. 短语}{迫切需要}
\ex{The slums are crying out for investment in clean water and schools.}
\ex{After the long march, the soldiers were crying out for a rest.}
\col{cry out（大声喊叫）}
\ex{She cried out in pain as the nurse removed the bandage.}
\col{cry out against（强烈抗议）}
\ex{Residents cried out against the new airport plan.}

\vw{cycle}{/ˈsaɪkl/}{n.\ \&\ vi.}

\sense{n.}{循环}
\col{a vicious cycle / break the cycle / the life cycle}
\ex{The region is trapped in a vicious cycle of poverty and poor education.}

\sense{n.}{自行车，摩托车}
\ex{We hired a couple of cycles and rode along the river bank.}

\sense{vi.}{骑自行车}
\ex{He cycles to school every day, rain or shine.}

\begin{gram}[-cycle 词族]
\textbf{cycle} $\to$ \textbf{recycle}（再循环）、\textbf{bicycle / tricycle / unicycle}（前缀 bi-, tri-, uni- 表数）、\textbf{cyclist}（骑车人）、\textbf{cyclone}（气旋“旋风”）。构词知识可用于推测生词词义，这是阅读理解词汇题的常设考点。
\end{gram}

\begin{origin}
希腊语 \textit{kyklos}“\textbf{圆环、车轮}”（原始印欧语 *\textit{kwel-}“\textbf{旋转、绕行}”）。这个“转圈”词根可能是英语中产出最多的家族之一：日耳曼一支给出 \textbf{wheel}（轮）与 \textbf{wel-}（\textit{waltz} 华尔兹“旋转之舞”）；希腊一支给出 \textbf{cycle}、\textbf{encyclopedia}（\textit{en-}+\textit{kyklos}+\textit{paideia}“圈内的教育”$\to$百科全书）、\textbf{cyclone}、\textbf{chakra}（梵语借词“轮”）；拉丁一支给出 \textbf{colony}（<*\textit{kwo-lo-}“绕圈耕作的农场”）与 \textbf{culture} 的一部分争论中的联系。一词根，三种文明，仍在转动。
\end{origin}

\end{document}
